\cotefile{127}
\scanpage{14}
% archive sheet number: 13
% L1
Considérons p.ex. \deleted{\unsure{de}} le faisceau d'anneaux
% L2
\unsure{standard|standart|structural} $O$ sur $C$, et considérons un
% L3
module $\deleted{\unsure{A}}F$ \inserted{quasi-cohérent} sur $O$. \deleted{En \illegible{2} $\unsure{\mathcal E}, F_{\unsure{\mathcal E}}$}
% L4
\deleted{\illegible{2}} Alors $F$ induit sur chaque
% L5
% slanted
\margin{à vérifier}
$\mathcal M_{/T}$ (pour
\diagram{\begin{tikzcd}[ampersand replacement=\&,row sep=0.15em,column sep=small]
X \& \& {}\\
{} \& T_0\arrow[ul,"\unsure{h|h_0}"']\arrow[r] \& T
\end{tikzcd}}) un faisceau \uline{quasi-cohérent},
% L6
avec \unsure{contravariance|isomorphismes} en l'objet variable (\diagram{\begin{tikzcd}[ampersand replacement=\&,row sep=0.18em,column sep=small]
X \& {}\\
T_0\arrow[u]\arrow[r] \& T
\end{tikzcd}})
% L7
de $C$, et de façon \deleted{\illegible{1}}\uline{cartésienne} \ldots. \deleted{De plus,}
% L8
on voit que l'on obtient en particulier
% L9
un ``faisceau formel'' sur le
% L10
complété formel d'un \deleted{$X$ \unsure{du|de}} $Z$ le long de
% L11
$X$ [pour une \unsure{imm|imm.} $\delta\colon X\unsure{-|\to} Z$, dont
% L12
les voisinages \unsure{infin.} sont $\in\operatorname{Ob}\,\underline{\unsure{S|Z|\mathcal S|\mathcal C}}$] et on
% L13
obtient alors, par une vérification essentiellement
% L14
facile
% L15
% slanted
\margin{\unsure{préciser|précision} et à vérifier.}
\[
H^*(\unsure{e|C|\ell|\mathcal C}_{/\delta},F)\simeq H^*(\widehat Z,F\mid\widehat Z)
\]
% L16
où le deuxième \unsure{mb} est pris comme
% L17
cohomologie dans la catégorie des systèmes
% L18
projectifs adiques sur le complété formel
% L19
$\widehat Z$ de $Z$ le long de $X$ \ldots\,.

